\subsection{The full tube action and the strip subgroups}

\begin{claim}{B.30}{Proved}{Explicit outer words for the antipodally paired tube model}
For $g\ge2$, put $n=2g$ and label the equatorial holes by $j\in\Z/n$,
with $j$ glued by a tube to $j+g$. The following presentation and
substitutions compute the entire exterior action requested on scan 467:
\[
 \pi_g=\left\langle x_0,\ldots,x_{g-1},p_0,\ldots,p_{g-1}
 \ \middle|\ R=1\right\rangle,
 \qquad
 R=x_0\cdots x_{g-1}
   \prod_{j=0}^{g-1}(p_j^{-1}x_j^{-1}p_j).
\]
Extend the notation, with all subscripts read modulo $n$, by
\[
 \begin{gathered}
 x_{j+g}=p_j^{-1}x_j^{-1}p_j,\qquad p_{j+g}=p_j^{-1},\\
 u_j=(x_0\cdots x_j)^{-1},\quad u_{n-1}=1,\qquad
 q_j=u_j^{-1}p_j u_{j+g-1}.
 \end{gathered}
\]
Let $r=\varepsilon_0\varepsilon_1$ and $s=\varepsilon_1$ as in
B.20. Representatives of their outer automorphisms, and of $\tau$ and
$\rho_k$, are specified on the $2g$ displayed generators by
\[
\begin{array}{c|cc}
 &x_j&p_j\\ \hline
 r&x_{j+1}&p_{j+1}\\
 s&u_{-j}^{-1}x_{-j}u_{-j}&q_{-j}\\
 \tau&u_{j+g-1}^{-1}x_{j+g}^{-1}u_{j+g-1}&q_j^{-1}\\
 \rho_k&x_j&x_j^{\delta_{jk}}p_j.
\end{array}
\]
The map to $\operatorname{Out}(\pi_g)$ defined by these substitutions
is the geometric action of
$\Z^g\rtimes(D_{4g}\times C_2)$ of B.20. The finite factor acts by
\[
 r\rho_jr^{-1}=\rho_{j+1},\qquad
 s\rho_js^{-1}=\rho_{-j},\qquad
 \tau\rho_j\tau^{-1}=\rho_j^{-1},
\]
with indices in these three formulas taken modulo $g$.
\end{claim}
\begin{proof}
We specify the base paths so that the formulas can be used as words,
rather than only as unbased pictures. Let $N,S$ be the north and south
poles in the sphere with the $n$ discs removed. Draw disjoint meridian
arcs $d_j:N\to S$ through the gap immediately following hole $j$.
The loop about that hole, based at $N$, is
$x_j=d_{j-1}d_j^{-1}$. Collapse the maximal-tree edge $d_{n-1}$ and
identify $N,S$ by that edge. The remaining arcs become
$u_j=(x_0\cdots x_j)^{-1}$, because the product telescopes to
$d_{n-1}d_j^{-1}$. In particular $\prod_{j=0}^{n-1}x_j=1$.

Let $p_j$ run from $N$ to the northern side of hole $j$, through the
tube to hole $j+g$, and back to $N$ on its northern side. The opposite
traversal is $p_{j+g}=p_j^{-1}$. Its boundary identification reads
$x_{j+g}=p_j^{-1}x_j^{-1}p_j$. Van Kampen therefore gives exactly the
presentation in the claim. The parallel path based at $S$, on the
southern side of the tube, is
$q_j=d_j^{-1}p_jd_{j+g-1}$ after inserting the northern paths, which
becomes the displayed formula on collapsing the tree edge.
These identities can also be taken as a two-vertex cellular description:
the sphere faces give the telescoping boundary relation and each tube
face gives the indicated boundary conjugation.

Rotation about the polar axis sends hole $j$ to $j+1$, preserves $N$,
and sends $p_j$ to $p_{j+1}$. This proves the $r$ row. The half-turn
$s$ about the equatorial axis through hole $0$ exchanges $N,S$ and
sends hole $j$ to $-j$. The southern boundary loop at hole $k$ is
$d_k^{-1}x_kd_k$, and the southern tube path is $q_k$, giving the
$s$ row. The antipodal map exchanges $N,S$, sends $j$ to $j+g$,
and reverses orientation. It therefore inverts the southern boundary
loop and reverses the corresponding tube path. The boundary-loop
formula written with $u_{j+g-1}$ equals that written with $u_{j+g}$,
since $u_{k-1}=x_ku_k$. This proves the $\tau$ row. Finally a positive
full twist of tube $k$ inserts $x_k$ into the crossing $p_k$ and
leaves all other chosen paths unchanged. This proves the $\rho_k$ row.
Each row comes from an actual homeomorphism, so it is an automorphism
of the stated group and has the prescribed orientation sign.

For $s$ and $\tau$ the chosen path from $N$ to the image of $N$ is
$d_{n-1}$. We consequently get actual automorphisms only after this
base-path choice. For example their squares are
\[
 s^2=\operatorname{Inn}(x_0),\qquad
 \tau^2=\operatorname{Inn}(x_0\cdots x_{g-1}),
\]
whereas $r^{2g}=1$ as an actual automorphism. These formulas follow
by concatenating the base path with its image: $s(d_{n-1})=d_0^{-1}$
and $\tau(d_{n-1})=d_{g-1}^{-1}$. Thus the displayed group relation
$s^2=\tau^2=1$ holds in the outer group with precisely the stated
inner ambiguity. The geometric symmetries satisfy $srs=r^{-1}$ and
commute with the antipodal involution on the surface. They preserve
and permute the tube collars; a positive symmetry preserves twist
sign and the antipodal map reverses it. This proves every relation
of the asserted semidirect product in the outer group. Changing the
maximal-tree edge conjugates each based comparison by its corresponding
base-path word and leaves this exterior action unchanged.

The accompanying programs check all relator images for $2\le g\le8$;
$r,s$ carry $R$ to a conjugate of $R$, and $\tau$ carries it to a
conjugate of $R^{-1}$. They additionally check all the finite-factor
and twist-conjugacy relations in the surface group, with the actual
inner conjugators above. For instance $srsr=1$ and
$\tau r\tau r^{-1}=\operatorname{Inn}(x_0\cdots x_g)$.
The word checker uses the symmetrized relator: its largest common
piece has length one, so the length-$4g$ presentation satisfies
$C'(1/6)$ for $g\ge2$, and the Dehn reduction is exact. These finite
certificates supplement the cellular proof for arbitrary $g$.
\end{proof}

\begin{claim}{B.31}{Proved}{The eight symmetries of one tube}
The stabilizer of tube $T_0$ in $D_{4g}\times\langle\tau\rangle$ is
$\langle r^g,s,\tau\rangle\cong C_2^3$. There are coordinates
$(\theta,t)\in(\R/2\pi\Z)\times[-1,1]$ on that tube in which its
three generators act as
\[
 r^g(\theta,t)=(-\theta,-t),\qquad
 s(\theta,t)=(\theta+\pi,t),\qquad
 \tau(\theta,t)=(\theta+\pi,-t).
\]
These formulas agree with the outer action in B.30 and with the
manuscript's free antipodal extension, after the explicit coordinate
change $\theta_{\rm new}=\theta_{\rm old}+\pi t/2$.
\end{claim}
\begin{proof}
The stabilizer of the unordered pair $\{0,g\}$ under the polygon's
dihedral group is $\{1,r^g,s,r^gs\}$: a stabilizing rotation is either
$1$ or $r^g$, and every stabilizing reflection is one of these followed
by $s$. The independent central antipodal factor doubles it. In a
symmetric tube, $r^g$ interchanges its ends and reverses its angular
direction, while $s$ preserves both ends and rotates the cross section
by $\pi$. The displayed representatives commute, are involutions,
and have eight distinct combinations; the antipodal representative
is free. They extend the same boundary maps as the corresponding
sphere symmetries and hence give the desired tube extension.

The formula of scan 453, retained through scans 461--462, is
$\tau_{\rm old}(\theta,t)=(\theta+\pi+\pi t,-t)$.
Conjugating it by the stated change gives exactly
$(\theta+\pi,-t)$; this also transports both boundary identifications.
Thus it is an explicit isomorphism of the two actions. The full twist
is $(\theta,t)\mapsto(\theta+\pi+\pi t,t)$ in either coordinate.
Conjugation by the two positive generators preserves its class, while
conjugation by the negative generator inverts its class, agreeing with
B.30. All eight maps and their products are checked symbolically in
\texttt{code/verify\_tube\_local.py}.
\end{proof}

\begin{claim}{B.32}{Proved}{Relations for arbitrary strips and a uniform obstruction to generation}
Choose for each $c_j$ one simple transversal $d_j$, using a strip in
the complement of the other tubes, so that
$i(c_i,d_j)=\delta_{ij}$. Let $W_j$ be a regular neighbourhood of
$c_j\cup d_j$, a one-holed torus, and put
$A_j=T_{c_j}=\rho_j$, $B_j=T_{d_j}$. Its supported subgroup is the
central extension $\langle A_j,B_j\rangle\cong B_3$ of B.21, with
quotient $SL_2(\Z)$. The exact pairwise relations are
\[
 [A_i,A_j]=1,\quad [A_i,B_j]=1\ (i\ne j),\qquad
 A_jB_jA_j=B_jA_jB_j.
\]
For distinct $i,j$, the two supported subgroups commute precisely when
$i(d_i,d_j)=0$. For intersection one, $B_i,B_j$ obey the braid
relation and do not commute; for intersection at least two they freely
generate a free group of rank two. For every such choice of one strip
per tube and every $g\ge2$, these $g$ supported subgroups, together
with $H$, generate a proper subgroup of $\mathcal T_g^+$.
\end{claim}
\begin{proof}
Each $W_j$ has genus one and one boundary component because the two
core curves meet transversely once. The extension-by-identity map and
its faithful boundary-twist kernel are the explicit maps of B.21.
The disjointness conditions give the first two commutation relations,
and the calculation on its two free generators in B.21 proves the
braid relation. For any two essential simple curves $a,b$, their
twists commute exactly when they are disjoint: conjugation sends
$T_b$ to $T_{T_a(b)}$, and
$i(T_a(b),b)=i(a,b)^2$. Hence intersecting curves cannot commute.
If $i(a,b)=1$, their regular neighbourhood is a one-holed torus, so
the same braid calculation applies.

For $k=i(a,b)\ge2$ the freedom assertion has an elementary ping-pong
proof. On isotopy classes of essential curves set
\[
 X_a=\{x:i(x,a)<i(x,b)\},\qquad
 X_b=\{x:i(x,b)<i(x,a)\}.
\]
Both are nonempty and disjoint. Counting crossings in a twist collar
and cancelling only crossings already present outside it gives
\[
 i(T_a^m(x),b)\ge |m|\,i(x,a)k-i(x,b).
\]
For $m\ne0$ and $x\in X_b$ this is strictly greater than
$i(x,a)=i(T_a^m(x),a)$, so $T_a^m(X_b)\subset X_a$.
The symmetric inclusion holds for powers of $T_b$. The ping-pong
argument on a reduced alternating word proves that it is nontrivial:
apply it to the appropriate one of $X_a,X_b$, inserting one further
nonzero power from the opposite factor when the first and last
factors agree. Thus $\langle T_a,T_b\rangle\cong\Z*\Z$.
These statements compute the required pairwise relations for arbitrary
intersections of the selected strips; no disjoint-strip convention is
being used to conclude the general generation assertion.

For that assertion work on $V=H_1(X_g;\F_2)$ with its nonsingular
intersection pairing. The $2g$ classes $[c_j],[d_j]$ form a basis:
pairing a linear relation first with the $[c_j]$ kills all coefficients
of the $[d_j]$, and pairing next with the $[d_j]$ kills the rest.
There is a unique quadratic refinement $q:V\to\F_2$ with
$q([c_j])=q([d_j])=1$. Explicitly, for this ordered basis $(v_i)$,
\[
 q\Bigl(\sum_i a_iv_i\Bigr)
 =\sum_i a_i+\sum_{i<j}a_ia_j\langle v_i,v_j\rangle.
\]
It obeys $q(v+w)=q(v)+q(w)+\langle v,w\rangle$ by expansion.
A twist about a curve of class $v$ acts by
$x\mapsto x+\langle x,v\rangle v$; substitution gives
\[
 q(T_vx)-q(x)=\langle x,v\rangle\bigl(q(v)+1\bigr).
\]
Every generator $A_j,B_j$ therefore preserves $q$, without any
restriction on the intersections among the $d_j$.

Join $c_1$ and $c_2$ by an embedded arc in their connected complement
and form a band sum $e$. This is a simple curve of nonzero class
$[e]=[c_1]+[c_2]$, so it is nonseparating. Since the two cores are
disjoint, $q([e])=0$. Taking $x=[d_1]$ gives
$\langle x,[e]\rangle=1$, so $T_e$ fails to preserve $q$.
Consequently $T_e$ belongs to $\mathcal T_g^+$ and lies outside the
group generated by the strip subgroups. This proves properness for
all choices under discussion. The quadratic calculation has been
exhaustively checked for every possible matrix of transversal
intersections modulo two in genera two, three and four; the general
identity is the displayed algebraic proof.

The rank-one comparison is thus the specified central quotient
$B_3\to SL_2(\Z)$ and its action on the homology of $W_j$, together
with the above braid and intersection relations. The subgroup they
generate fixes the explicit quadratic refinement. This supplies the
precise algebraic maps and the explicit obstruction governing the
proposed generation in scan 470.
\end{proof}
